% a05_twoloop_iv: proof of lem:a05:D45 (iv), the |y_21| bound on eps Z^3 (input from a05_cnum_NEG.tex)
\paragraph{Proof of \cref{lem:a05:D45}(iv).} The lattice is $\eps\Z^3$ with the conventions of the paragraph ``Lattice formulas'': $s=M^2$, $C(q)=(\hat q^2+s)^{-1}$, $D(k)=(\hat k^2+\lmu^2)^{-1}$, $\langle\cdot\rangle$ the integral over $\BZ_\eps$, $\bar\BZ_\eps$ the closed zone, $\hat x_\mu\defeq\frac2\eps\sin\frac{\eps x_\mu}2$. Every integrand below is smooth and $\frac{2\pi}\eps\Z^3$-periodic in each integration variable, so $k\mapsto p-k$ is a permitted substitution for every real $p$. Put
\begin{gather*}
B_s(p)\defeq\langle D(k)C(p-k)\rangle_k,\qquad \Xi_s(p)\defeq\langle D(k)C(p-k)^2\rangle_k=-\partial_sB_s(p),\qquad J^\alpha_s(p)\defeq\langle D(k)C(p-k)\widehat{(p-k)}_\alpha^{\,2}\rangle_k,\\
s_0\defeq\langle D\rangle,\qquad b_s\defeq\langle C^2\rangle,\qquad F_s(p)\defeq\langle\widehat{(p+r)}^2D(p-r)C(r)\rangle_r,\qquad G_s(p)\defeq\partial_sF_s(p)=-\langle\widehat{(p+r)}^2D(p-r)C(r)^2\rangle_r .
\end{gather*}
$F_s$ is the photon-exchange graph and $s_0\sum_\mu\cos\eps p_\mu$ the seagull tadpole at external momentum $p$: the function $\hat\Pi$ of \cref{lem:a05:deriv}, at spacing $\eps$ and after the substitution $r=p+k$, is $s_0\sum_\mu\cos\eps p_\mu-F_s(p)$. Hence the self-energy per real component is
\begin{equation}\label{eq:a05:sigeps}
\sigma_\eps(p)=-\bar\pi_\eps(p)=F_s(p)-F_s(0)+s_0\sum_\mu(1-\cos\eps p_\mu)=F_s(p)-F_s(0)+\tfrac12\eps^2s_0\,\hat p^2 ,
\end{equation}
and since $s_0$ does not depend on $s$, $\partial_s\sigma_\eps(q)=G_s(q)-G_s(0)$. Differentiating (iv) of the lattice formulas,
\begin{equation}\label{eq:a05:y21}
y_{21}=\partial_sd_{21}=-16\big\langle-2C^3\sigma_\eps+C^2\,\partial_s\sigma_\eps\big\rangle .
\end{equation}

\begin{lemma}[elementary lattice inequalities, continued]\label{lem:a05:el2}
Let $u,v\in\R^3$, $\tilde a\defeq\frac\pi2\sqrt s$, $\tilde b\defeq\frac\pi2\lmu$, $\tilde G_a(w)\defeq(|w|^2+a^2)^{-1}$ on $\R^3$.
\begin{enumerate}[label=\textup{(\alph*)},nosep,leftmargin=*]
\item $|\hat u_\mu-\hat v_\mu|\le|u_\mu-v_\mu|$, $|\widehat{(u+v)}_\mu|\le|\hat u_\mu|+|\hat v_\mu|$, $|\widehat{(2v)}_\mu|\le2|\hat v_\mu|$; hence $\widehat{(u+v)}^2\le2\,\widehat{(u-v)}^2+8\hat v^2$.
\item With $E(v)\defeq\hat v^2$: $\partial_\mu E=\frac2\eps\sin\eps v_\mu$, $|\partial_\mu E|\le2|\hat v_\mu|$, $\partial_\mu\partial_\nu E=2\cos(\eps v_\mu)\delta_{\mu\nu}$; and $|w\cdot\operatorname{Hess}C(v)\,w|\le10\,C(v)^2|w|^2$ for all $w\in\R^3$.
\item For $p\in\bar\BZ_\eps$ and $k\in\BZ_\eps$: $C(p-k)\le\frac{\pi^2}4\sum_{n\in\{-1,0,1\}^3}\tilde G_{\tilde a}\big(p-k-\tfrac{2\pi}\eps n\big)$ and $C(p-k)^2\le\big(\frac{\pi^2}4\big)^2\sum_n\tilde G_{\tilde a}\big(p-k-\tfrac{2\pi}\eps n\big)^2$; for $k\in\BZ_\eps$, $C(k)\le\frac{\pi^2}4\tilde G_{\tilde a}(k)$ and $D(k)\le\frac{\pi^2}4\tilde G_{\tilde b}(k)$.
\item $s_0\le d_0/\eps$ with $d_0=\intk\hat k^{-2}\le\frac{\sqrt3\,\pi}8$ ($d_0$ as in \cref{lem:a05:J}); $0\le J^\alpha_s(p)\le s_0$.
\item $\int_{\R^3}\tilde G_b(k)\tilde G_a(P-k)^2\frac{d^3k}{(2\pi)^3}=\dfrac1{8\pi a\,((a+b)^2+|P|^2)}\le\dfrac1{8\pi a(a+b)^2}$.
\end{enumerate}
\end{lemma}
\begin{proof}
(a) The derivative of $x\mapsto\frac2\eps\sin\frac{\eps x}2$ is $\cos\frac{\eps x}2$, of modulus $\le1$; the addition theorem gives the second claim and $|\sin2y|\le2|\sin y|$ the third. Then $|\widehat{(u+v)}_\mu|=|\widehat{((u-v)+2v)}_\mu|\le|\widehat{(u-v)}_\mu|+2|\hat v_\mu|$ and $(x+y)^2\le2x^2+2y^2$; sum over $\mu$.
(b) $|\frac2\eps\sin\eps v_\mu|=\frac4\eps|\sin\frac{\eps v_\mu}2\cos\frac{\eps v_\mu}2|\le2|\hat v_\mu|$. With $\partial_\mu C=-C^2\partial_\mu E$: $w\cdot\operatorname{Hess}C\,w=2C^3(w\cdot\nabla E)^2-2C^2\sum_\mu\cos(\eps v_\mu)w_\mu^2$, and $|w\cdot\nabla E|\le\sum_\mu2|w_\mu||\hat v_\mu|\le2|w||\hat v|$, so the modulus is $\le(8C^3\hat v^2+2C^2)|w|^2\le10C^2|w|^2$, since $C\hat v^2\le1$.
(c) The components of $p-k$ lie in $[-\frac{2\pi}\eps,\frac{2\pi}\eps]$, so a nearest point $\bar n$ of $\frac{2\pi}\eps\Z^3$ lies in $\frac{2\pi}\eps\{-1,0,1\}^3$. By \cref{lem:a05:el}(c) at spacing $\eps$, $\widehat{(p-k)}^2\ge\frac4{\pi^2}|p-k-\bar n|^2$, so $C(p-k)\le(\frac4{\pi^2}|p-k-\bar n|^2+s)^{-1}=\frac{\pi^2}4\tilde G_{\tilde a}(p-k-\bar n)$; a maximum of positive terms is at most their sum. For $k\in\BZ_\eps$, $\bar n=0$.
(d) $k=k'/\eps$ gives $\langle\hat k^{-2}\rangle=d_0/\eps$, and $d_0\le\frac{\pi^2}4\int_{|k|\le\sqrt3\pi}\frac{d^3k}{(2\pi)^3|k|^2}=\frac{\pi^2}4\cdot\frac{\sqrt3}{2\pi}$. $J^\alpha_s\le s_0$ because $C(v)\hat v_\alpha^2\le C(v)\hat v^2\le1$.
(e) By \cref{lem:a05:bub}(i) the integral with $\tilde G_a$ in place of $\tilde G_a^2$ is $\frac{\arctan(|P|/(a+b))}{4\pi|P|}$; apply $-\partial/\partial(a^2)=-(2a)^{-1}\partial_a$ and $\partial_a\arctan\frac x{a+b}=-\frac x{(a+b)^2+x^2}$.
\end{proof}

\begin{lemma}[the decomposition on $\eps\Z^3$]\label{lem:a05:Kdec}
For every $p\in\R^3$,
\[
F_s(p)-F_s(0)=2\hat p^2B_s(p)+(\lmu^2-2s)\big(B_s(p)-B_s(0)\big)-\eps^2\sum_\alpha\hat p_\alpha^2J^\alpha_s(p),
\]
hence $\sigma_\eps(p)=2\hat p^2B_s(p)+(\lmu^2-2s)(B_s(p)-B_s(0))+2\rho_\eps(p)$ with $\rho_\eps(p)\defeq-\frac12\eps^2\sum_\alpha\hat p_\alpha^2J^\alpha_s(p)+\frac14\eps^2s_0\hat p^2$. (This is \cref{lem:a05:dec} at spacing $\eps$.)
\end{lemma}
\begin{proof}
\Cref{lem:a05:trig} with $u=\eps x$, $w=\eps y$, divided by $\eps^2$, reads $\widehat{(x+y)}^2=2\hat x^2+2\hat y^2-\eps^2\hat x^2\hat y^2-\widehat{(x-y)}^2$ componentwise. With $x=p_\alpha$, $y=(p-k)_\alpha$, summed over $\alpha$: $\widehat{(2p-k)}^2=2\hat p^2+2\widehat{(p-k)}^2-\eps^2\sum_\alpha\hat p_\alpha^2\widehat{(p-k)}_\alpha^2-\hat k^2$. Substitute $r=p-k$ in $F_s$ and insert $\widehat{(p-k)}^2=C(p-k)^{-1}-s$, $\hat k^2=D(k)^{-1}-\lmu^2$:
\[
F_s(p)=\big\langle D(k)C(p-k)\big[2\hat p^2-2s+\lmu^2\big]\big\rangle_k+2s_0-\langle C\rangle-\eps^2\sum_\alpha\hat p_\alpha^2J^\alpha_s(p),
\]
using $\langle C(p-k)\rangle_k=\langle C\rangle$. Subtract the value at $p=0$, where $\hat p=0$: the linearly divergent constants $2s_0-\langle C\rangle$ cancel exactly. The second form follows from \eqref{eq:a05:sigeps}.
\end{proof}

\begin{lemma}[one-loop blocks with explicit $s$]\label{lem:a05:P}
For every $\eps>0$ and $s>0$:
\begin{enumerate}[label=\textup{(P\arabic*)},nosep,leftmargin=*]
\item $b_s\le\dfrac{\pi^2}{64\sqrt s}$;
\item $B_s(0)\le\dfrac{\pi^2}{32(\sqrt s+\lmu)}$ and $\sup_{p\in\bar\BZ_\eps}B_s(p)\le\dfrac{27\pi^2}{32(\sqrt s+\lmu)}$;
\item $\sup_{p\in\bar\BZ_\eps}\Xi_s(p)\le\dfrac{27\pi^2}{64\,s^{3/2}}$;
\item $|G_s(p)|\le2b_s+8B_s(p)$ for $p\in\bar\BZ_\eps$; hence $|\partial_s\sigma_\eps(q)|\le4b_s+16\sup_{\bar\BZ_\eps}B_s$ for $q\in\BZ_\eps$;
\item $|\sigma_\eps(p)|\le2\hat p^2B_s(p)+5(\lmu^2+2s)\,|p|^2\sup_{\bar\BZ_\eps}\Xi_s+\frac32\eps^2s_0\hat p^2$ for $p\in\BZ_\eps$.
\end{enumerate}
\end{lemma}
\begin{proof}
(P1) By \cref{lem:a05:el2}(c), $b_s\le(\frac{\pi^2}4)^2\int_{\R^3}\tilde G_{\tilde a}^2\frac{d^3q}{(2\pi)^3}=\frac{\pi^4}{16}\cdot\frac1{8\pi\tilde a}=\frac{\pi^2}{64\sqrt s}$.
(P2) At $p=0$ both factors are at $k\in\BZ_\eps$: $B_s(0)\le(\frac{\pi^2}4)^2\int\tilde G_{\tilde b}\tilde G_{\tilde a}=\frac{\pi^4}{16}\cdot\frac1{4\pi(\tilde a+\tilde b)}$ by \cref{lem:a05:bub}(i). At general $p\in\bar\BZ_\eps$, \cref{lem:a05:el2}(c) supplies $27$ images, each integral being $\le\frac1{4\pi(\tilde a+\tilde b)}$.
(P3) \Cref{lem:a05:el2}(c), (e): $\Xi_s(p)\le(\frac{\pi^2}4)^3\sum_n\int\tilde G_{\tilde b}(k)\tilde G_{\tilde a}(p-k-\tfrac{2\pi}\eps n)^2\frac{d^3k}{(2\pi)^3}\le\frac{27(\pi^2/4)^3}{8\pi\tilde a(\tilde a+\tilde b)^2}\le\frac{27\pi^6/64}{8\pi(\pi/2)^3s^{3/2}}=\frac{27\pi^2}{64\,s^{3/2}}$.
(P4) \Cref{lem:a05:el2}(a) with $u=p$, $v=r$: $\widehat{(p+r)}^2D(p-r)C(r)^2\le\big[2\widehat{(p-r)}^2D(p-r)+8\hat r^2D(p-r)\big]C(r)^2\le2C(r)^2+8D(p-r)C(r)$, by $\hat x^2(\hat x^2+\lmu^2)^{-1}\le1$ and $\hat r^2C(r)\le1$. Integrate over $r$; $\langle D(p-r)C(r)\rangle_r=B_s(p)$. The second claim is $\partial_s\sigma_\eps=G_s(q)-G_s(0)$.
(P5) In \cref{lem:a05:Kdec}, $\eps^2\sum_\alpha\hat p_\alpha^2J^\alpha_s\le\eps^2\hat p^2s_0$ by \cref{lem:a05:el2}(d), and the seagull term is $\frac12\eps^2s_0\hat p^2$. For the middle term: $B_s$ is smooth on $\R^3$ and even ($k\mapsto-k$; $D$, $C$ even), so $\nabla B_s(0)=0$ and $B_s(p)-B_s(0)=\int_0^1(1-t)\,(p\cdot\nabla)^2B_s(tp)\,dt$ with $(p\cdot\nabla)^2B_s(tp)=\langle D(k)\,p\cdot\operatorname{Hess}C(tp-k)\,p\rangle_k$. By \cref{lem:a05:el2}(b), $|(p\cdot\nabla)^2B_s(tp)|\le10|p|^2\Xi_s(tp)$ with $tp\in\bar\BZ_\eps$, and $\int_0^1(1-t)dt=\frac12$. Finally $|\lmu^2-2s|\le\lmu^2+2s$.
\end{proof}

\begin{proof}[Proof of (iv)]
By \eqref{eq:a05:y21}, $|y_{21}|\le16\big[2\langle C^3|\sigma_\eps|\rangle+\langle C^2|\partial_s\sigma_\eps|\rangle\big]$. For $q\in\BZ_\eps$, $C^3\hat q^2\le C^2$ and $C^3|q|^2\le\frac{\pi^2}4C^3\hat q^2\le\frac{\pi^2}4C^2$ (\cref{lem:a05:el}(c)). So \cref{lem:a05:P}(P5) gives $2\langle C^3|\sigma_\eps|\rangle\le b_s\big[4\sup B_s+\frac{5\pi^2}2(\lmu^2+2s)\sup\Xi_s+3\eps^2s_0\big]$, and (P4) gives $\langle C^2|\partial_s\sigma_\eps|\rangle\le b_s\big[4b_s+16\sup B_s\big]$, the suprema over $\bar\BZ_\eps$. Hence
\[
|y_{21}|\le16\,b_s\Big[20\sup B_s+\frac{5\pi^2}2(\lmu^2+2s)\sup\Xi_s+4b_s+3\eps^2s_0\Big].
\]
Insert (P1)--(P3) and \cref{lem:a05:el2}(d), dropping $\lmu\ge0$ from the denominators: $16b_s\le\frac{\pi^2}{4\sqrt s}$; $20\sup B_s\le\frac{135}8\frac{\pi^2}{\sqrt s}$; $\frac{5\pi^2}2(\lmu^2+2s)\cdot\frac{27\pi^2}{64s^{3/2}}=\frac{135\pi^4}{128}\big(2+\frac{\lmu^2}s\big)\frac1{\sqrt s}$; $4b_s\le\frac{\pi^2}{16\sqrt s}$; $3\eps^2s_0\le\frac{3\sqrt3\,\pi}8\eps$. Collecting,
\[
|y_{21}(s;\eps)|\le\frac{\pi^2}4\Big[\frac{135}8\pi^2+\frac{135\pi^4}{128}\Big(2+\frac{\lmu^2}s\Big)+\frac{\pi^2}{16}\Big]\frac1s+\frac{\pi^2}4\cdot\frac{3\sqrt3\,\pi}8\cdot\frac\eps{\sqrt s}=\frac{C_{21}(s)}s+C'_{21}\frac\eps{\sqrt s},
\]
and $C_{21}$ is non-increasing in $s$ because $\lmu^2/s$ is. The only coefficient that grows with $s$, $|\lmu^2-2s|$, has been paired with $B_s(p)-B_s(0)=O(|p|^2s^{-3/2})$, the Taylor step being taken on the scalar line.
\end{proof}
